\chapter{Residual Diagnostics and Censoring}\label{ch10:residuals}

\section{The residual, defined narrowly}\label{ch10:sec:definition}

For an admitted active-region prediction, the residual is the signed
difference between what was observed and what the calibrated model
predicted:
\[
\delta_y \;=\; y_{\mathrm{obs}} - y_{\mathrm{pred}},
\]
where $y_{\mathrm{pred}} = a / r_{\mathrm{eff}}$ is the registered
prediction from a declared alpha dose $a$ and a calibrated effective
resistance $r_{\mathrm{eff}}$, and $y_{\mathrm{obs}}$ is an instrumented
observation of the flow magnitude at the same test point. In the
implementation the emitted quantity is named \texttt{delta\_I} (with the
observation and prediction in \textsc{activation} units and the residual
of quantity kind \texttt{FLOW\_RESIDUAL}); this chapter writes
$\delta_y$ for continuity with the response-coordinate notation of the
seed paper, and nothing turns on the letter. Alongside the residual, the
implementation emits a relative error
$\varepsilon = |\delta_y| / \max(y_{\mathrm{pred}}, \phi)$, where $\phi$
is the profile's declared residual floor --- a strictly positive
tolerance that prevents division by predictions compatible with zero ---
and, in the uncensored case only, an apparent resistance
$a / y_{\mathrm{obs}}$ (explicitly role-tagged
\texttt{APPARENT\_SAME\_SAMPLE}, the role the prediction gate refuses as
input) and a resistance delta against the calibrated value. Every one of
these quantities carries propagated uncertainty; a residual without an
uncertainty interval is unrepresentable, since the result type rejects
construction when either the residual or the relative error lacks one.

That is the entire positive content of the layer. The residual is a
number with an interval, computed by subtraction, bound to its
derivation. Everything else in this chapter concerns what the candidate
refuses to let that number mean.

\section{Admission requirements}\label{ch10:sec:admission}

The diagnostic entry point is gated more heavily than any other
operation in the candidate, because its output is the one most liable to
overinterpretation (Section~\ref{ch10:sec:tensions}). The requirements,
as implemented, are these.

\begin{ruleenv}[Residual admission]\label{ch10:rule:admission}
A residual is computed only when all of the following hold.
\begin{enumerate}
\item \emph{Instrumented, held-out observation.} The observed value is of
quantity kind \texttt{FLOW\_MAGNITUDE} in \textsc{activation} units and
of measurement kind \texttt{INSTRUMENTED} --- which itself requires an
instrument identity, a test-point identity, a calibration-profile
binding, and raw artifact references --- with uncertainty, a recorded
sample size, linked controls, and a dataset-split reference equal to the
profile's held-out validation split. A score, a behavioral proxy, or a
model-derived value cannot occupy the observation slot, and a
calibration-split observation cannot masquerade as held-out.
\item \emph{Matching test points.} Drive, prediction, observation, and
calibrated resistance must bind one common test-point identity; any
disagreement, or any missing test point, is a validation error.
\item \emph{Profile identity.} The operating assessment must bind the
context's exact calibration-profile digest and a test point registered
in that profile; the context itself must have admitted the profile as
reviewed, drift-current, and valid at the caller-supplied replay time.
\item \emph{Registered prediction.} The predicted value is not accepted
as a number: the gate recomputes the prediction from the presented drive
and calibrated resistance through the full prediction pathway --- with
all of that pathway's own gates: design-value drive within the
calibrated alpha range, calibrated-held-out or replayed
network-equivalent resistance bound to the profile's fit method, fit
receipt, source and control evidence --- and requires digest equality
between the recomputation and the presented prediction. A prediction
that did not come from the registered pathway cannot have a residual.
\item \emph{Calibrated resistance.} The resistance carries positive
uncertainty bounds, a sample size, and an admissible role; apparent
same-sample resistance is rejected at the prediction gate, which
forecloses the circular construction $r = a/y$, $\hat y = a/r$,
$\delta_y = 0$, ``validated.''
\item \emph{Active region, assessed outcome.} The operating assessment's
region must be \texttt{ACTIVE} and its semantic outcome must not be
\texttt{UNASSESSED}; numeric assessments are replayed against their
bound sweep evidence before use. The no-diagnosis conditions of
Chapter~\ref{ch07:classification} apply in full.
\item \emph{Propagated uncertainty.} Prediction and observation must
both carry uncertainty; the residual interval is computed by interval
subtraction (observed lower minus predicted upper, observed upper minus
predicted lower), the standard uncertainty in quadrature, and the
confidence level as the minimum of the inputs. The relative error's
denominators are floored at $\phi$ throughout.
\end{enumerate}
\end{ruleenv}

The cumulative effect is that a residual exists only inside a fully
reconstructed derivation: same profile, same test point, held-out data,
registered prediction, calibrated parameter, replayable assessment. There
is no API by which two bare numbers can be subtracted into a
\texttt{FLOW\_RESIDUAL}.

Requirement 4 deserves one further paragraph, because it is the
admission rule with the longest reach. The gate does not compare the
presented prediction to a recomputation \emph{of its value}; it compares
canonical digests, which cover value, unit, uncertainty, sample size,
provenance, and operation binding jointly. A presented prediction with
the right value but the wrong derivation --- computed from a different
resistance, under a different profile revision, or with its uncertainty
altered --- fails the digest comparison exactly as a wrong value would.
This is the anti-circularity rule of the seed paper made operational at
the diagnostic layer: the same-sample construction ($r = a/y$ on a
sample, then $\hat y = a/r$ on that sample, then $\delta_y = 0$,
``validated'') cannot reach the residual gate, because the apparent
same-sample resistance role is rejected upstream by the prediction
pathway the gate replays, and no alternative pathway into the
prediction slot exists. Verification of this closure --- that the
circular route fails at every seam --- is part of the bounded software
evidence; it is a property of the software, established by tests, in the
verification sense rather than the validation sense \cite{sargent}.

\begin{example}[Illustrative admitted residual]\label{ch10:ex:residual}
All numbers illustrative. Under one admitted context: drive $a = 4.0$
alpha (design value, inside the calibrated grid); calibrated resistance
$r_{\mathrm{eff}} = 2.0$ alpha-per-activation, interval $[1.8, 2.2]$,
standard uncertainty $0.1$, sample size $40$, role
calibrated-held-out with the profile's fit receipt as sole derivation
parent. The registered prediction is $y_{\mathrm{pred}} = 4.0 / 2.0 =
2.0$ activation units, with interval
$[\,4.0/2.2,\; 4.0/1.8\,] \approx [1.82, 2.22]$ and first-order standard
uncertainty $\approx 0.10$. An instrumented held-out observation at the
same test point reads $y_{\mathrm{obs}} = 2.6$ with interval
$[2.3, 2.9]$ and standard uncertainty $0.15$. The residual is
$\delta_y = 0.6$ with interval
$[\,2.3 - 2.22,\; 2.9 - 1.82\,] \approx [0.08, 1.08]$ and standard
uncertainty $\sqrt{0.15^2 + 0.10^2} \approx 0.18$; the relative error is
$0.6 / 2.0 = 0.30$; the apparent resistance is $4.0 / 2.6 \approx 1.54$
(role \texttt{APPARENT\_SAME\_SAMPLE}); the resistance delta is
$\approx -0.46$ with its own propagated interval. The residual interval
excludes zero, so the record is uncensored and carries all four
quantities. What may be said: the calibrated account underpredicts this
observation by $0.6$, more than its combined uncertainty. What may not
be said: why. Every mechanism in the list of
Section~\ref{ch10:sec:nonident} remains open, including the boring
ones.
\end{example}

\begin{example}[Illustrative censored case]\label{ch10:ex:censored}
All numbers illustrative. Same context and prediction, but the
observation reads $y_{\mathrm{obs}} = 0.1$ with interval $[-0.05, 0.25]$
(zero-compatible: the instrument cannot distinguish this reading from
silence at the stated confidence). The residual $\delta_y = -1.9$ with
interval $[\,-0.05 - 2.22,\; 0.25 - 1.82\,] \approx [-2.27, -1.57]$ is
emitted --- it is informative, and its interval excludes zero --- but
the apparent resistance $4.0 / 0.1 = 40$ is \emph{not} emitted, because
the denominator's interval includes zero and the value $40$ is an
artifact of where the point estimate fell: had the point estimate been
$0.01$, the ``apparent resistance'' would have been $400$ from data the
instrument cannot distinguish from the same silence. The result carries
the censoring marker \texttt{OBSERVED\_INTERVAL\_INCLUDES\_ZERO} and its
recorded reason, and the censored result type structurally cannot carry
resistance estimates. The flat-response ambiguity of
Chapter~\ref{ch07:classification} --- cutoff, saturation, wrong
direction, insensitivity, measurement failure --- applies to this record
in full, which is exactly why no ratio is allowed to lend it spurious
precision.
\end{example}

\section{Censoring at zero-compatible response}\label{ch10:sec:censoring}

\begin{definition}[Censoring rule]\label{ch10:def:censoring}
If the observed value's uncertainty interval includes or touches zero
(lower bound $\le 0$), the diagnostic emits the residual and relative
error --- both remain well-defined differences --- but \emph{censors} the
ratio quantities: no apparent resistance and no resistance delta are
constructed, the result carries the censoring marker
\texttt{OBSERVED\_INTERVAL\_INCLUDES\_ZERO} with a recorded reason, and
the censored result type rejects any attempt to attach resistance
estimates to it. Conversely an uncensored result must carry both ratio
quantities and no reason; the two states are mutually exclusive by
construction.
\end{definition}

The statistical rationale is elementary but worth stating carefully,
because the source material violates it. The apparent resistance is the
ratio $a / y_{\mathrm{obs}}$. A ratio whose denominator is compatible
with zero has no finite moments and no usable interval: as the
denominator's interval approaches zero from above, the ratio's upper
bound diverges, and once the interval includes zero the ratio's sampling
distribution places mass at arbitrarily large values of both signs. Any
finite number reported in that situation is an artifact of where the
point estimate happened to fall, not a summary of what the data support.
The honest treatments are the standard ones for censored data: report
the censoring event itself as the datum (the observation is
indistinguishable from zero at the stated confidence), carry the bound
rather than a point value, and refuse downstream arithmetic that would
consume the undefined ratio. The candidate implements the refusal
structurally --- the censored result type \emph{cannot} carry resistance
estimates --- rather than advisorily.

Prediction-side degeneracy is handled by the complementary mechanism:
denominators involving the predicted value are floored at the profile's
residual floor $\phi > 0$, so a prediction at or near zero yields a
relative error computed against $\phi$ rather than an unbounded ratio,
and the floor value itself is a declared, digested profile tolerance
rather than a hidden constant. Between the censoring rule and the floor,
there is no code path on which a division by a zero-compatible quantity
produces a reported diagnostic.

Note also what censoring does \emph{not} do: it does not suppress the
residual. A zero-compatible observation against a nonzero prediction is
informative --- it is the flat-response situation of
Chapter~\ref{ch07:classification}, with its five-way causal ambiguity
--- and the signed difference with its interval is retained. What is
suppressed is exactly the derived quantity whose statistical meaning has
evaporated.

\section{Residuals are non-identifying}\label{ch10:sec:nonident}

\subsection{The argument}

The central claim of this chapter, and the one the source audit shows
the source material itself half-knows, is that a residual --- however
clean its derivation --- cannot identify its own cause. The argument
deserves to be made in full rather than gestured at.

A residual is a one-dimensional summary: a signed magnitude (with
interval) of disagreement between one observation and one model
prediction at one point. Consider the space of generating situations
compatible with a given positive residual $\delta_y > 0$ (observed flow
exceeding prediction; the negative case is symmetric). At minimum, the
following mechanisms each produce it:

\begin{enumerate}
\item \emph{Hidden facilitation or policy.} Some mechanism not in the
model assists the driven response --- the ``hidden bypass'' reading: the
effective opposition is lower than calibrated.
\item \emph{Bypass of a modeled constraint.} A constraint that the
calibrated resistance implicitly includes is inoperative in this
context, for reasons benign or otherwise.
\item \emph{Opposition elsewhere.} A mechanism opposing a
\emph{different} pathway redirects activation such that the measured
projection increases; the causal action is not on the measured path at
all.
\item \emph{Measurement error.} The instrument chain --- hooks,
projections, normalization, full-scale binding --- inflates the
observation; nothing in the model is wrong.
\item \emph{Model misspecification.} The true local response is
nonlinear, and the linear prediction undershoots at this dose; the
discrepancy is structural, not mechanistic \cite{kennedyohagan}.
\item \emph{Calibration transfer failure.} The resistance was fit under
conditions (corpus, seeds, drift state) that no longer obtain, and the
parameter, not the mechanism, changed.
\item \emph{Sampling fluctuation.} The residual's own interval, honestly
computed, may include values consistent with all of the above at once.
\end{enumerate}

These mechanisms are observationally equivalent at the level of the
single residual: the datum is one number, the hypothesis space is many-dimensional,
and the map from mechanisms to residual sign is many-to-one in both
directions (several mechanisms produce either sign, depending on
parameters). Sign alone therefore cannot select among them --- not
weakly, not suggestively, but as a matter of the arithmetic of
identification.

The point can be sharpened in the vocabulary of causal inference. Call
two generating situations \emph{residual-equivalent} if they induce the
same distribution over $(a, y_{\mathrm{obs}})$ at the observed test
points. Mechanism attribution from residuals is well-posed only if the
attribution target is constant on residual-equivalence classes --- that
is, only if the observational distribution identifies the mechanism.
The mechanism list above exhibits, constructively, residual-equivalent
situations with different mechanisms: a hidden facilitation of the
measured path and a measurement-chain gain error can be tuned to
produce identical $(a, y_{\mathrm{obs}})$ distributions, as can a
bypassed constraint and a miscalibrated $r_{\mathrm{eff}}$. Hence the
attribution target is not constant on the equivalence classes, and no
amount of statistical sophistication applied to the same observational
data can recover it; what changes the situation is intervention. This
is the standard identification requirement: causal quantities are
defined relative to an explicit causal model and are computable from
data only when the model plus the data identify them, typically via
interventions or assumptions that break the observational ties
\cite{pearl}. The candidate supplies no causal graph over the mechanism
space, asserts no identifying assumptions, and admits no
do-operations; its residuals are observational discrepancies, and it
says so in their method identity. Distinguishing the mechanisms
requires interventions that separate them: manipulating the suspected
constraint while holding the dose fixed, ablating the suspected bypass,
permuting the labels the instrument depends on, replicating under
re-calibration --- each an experiment with its own controls, several of
them present in the registered closure list. The conclusion is not that
residual-based mechanism attribution is difficult; it is that, from the
residual alone, it is undefined.

There is also a model-side face of the same coin, familiar from the
calibration literature: an observed discrepancy decomposes, in
principle, into parameter error, structural model discrepancy,
observation error, and numerical error, and the decomposition is itself
unidentified without prior structure --- fitting the parameter to absorb
the discrepancy simply relocates the problem \cite{kennedyohagan}.
Mechanism 5 and mechanism 6 in the list are this decomposition's
components wearing mechanistic clothes; the candidate's refusal to
update $r_{\mathrm{eff}}$ from residuals (no feedback path exists from
the diagnostic layer to the profile) is the structural form of refusing
to let the parameter absorb what may be discrepancy.

\begin{example}[One residual, three stories]\label{ch10:ex:stories}
All numbers illustrative. Take the uncensored record of
Example~\ref{ch10:ex:residual}: $\delta_y = +0.6$, interval
$[0.08, 1.08]$. Story A (hidden facilitation): a mechanism outside the
model contributes response along the measured direction; true
opposition is lower than calibrated. Story B (measurement): the
full-scale value bound at this test point was set from a corpus
percentile that undersamples high-activation contexts, inflating the
instrumented reading by a factor compatible with $+0.6$. Story C
(misspecification): the true response is mildly convex in $a$ over the
grid, the linear fit undershoots at $a = 4.0$, and the same board would
show a \emph{negative} residual at $a = 1.0$. All three stories
reproduce the record; the second and third involve no hidden mechanism
at all; and the third makes a testable prediction (sign reversal at low
dose) that the first two do not --- which is a statement about what to
probe next, not about which story is true. This is the entire licensed
use of the record.
\end{example}

\begin{ruleenv}[Non-identification]\label{ch10:rule:nonident}
No implemented output of the residual layer asserts, encodes, or ranks a
generating mechanism. The emitted quantities are the residual, the
relative error, and (uncensored only) the apparent resistance and
resistance delta --- each a descriptive statistic with propagated
uncertainty, none a diagnosis. The method identity stamped on every
residual's provenance names the contract:
\texttt{CAS-OHM-V1-NONIDENTIFYING-RESIDUAL}.
\end{ruleenv}

\subsection{What residuals are good for}\label{ch10:sec:goodfor}

Non-identification does not make residuals useless; it makes them
\emph{ordinal evidence for allocation of further measurement}. A large
residual (relative to its interval) at one test point, against a small
one at another, is a defensible reason to spend the next intervention
budget at the first point --- not because the residual says what is
there, but because it says the calibrated model's account is worse
there. Ranking follow-up probes is a decision about experiments, not a
claim about mechanisms, and it is robust to the non-identification
argument precisely because every mechanism on the list above, including
boring measurement error, is something a follow-up probe can help
separate. The resistance delta serves the same triage role in parameter
space: a delta whose interval excludes zero flags the calibration, the
transfer conditions, or the linearity assumption for re-examination,
without saying which.

The ranking criterion itself can be made precise without exceeding the
license. For admitted, uncensored residual records $j = 1, \ldots, k$
over one profile, a defensible ordering statistic is the
uncertainty-scaled magnitude $|\delta_y^{(j)}| / s^{(j)}$, where
$s^{(j)}$ is the record's propagated standard uncertainty --- large
values mark points where the disagreement is large relative to what the
error budget explains. The statistic inherits every caveat of its
inputs: it compares only within one profile (cross-profile residuals
are in incomparable coordinates), only among uncensored records
(censored records enter the queue by their censoring event, not by a
ratio), and only as a priority order, since the statistic's
distribution under the no-mechanism null is exactly what the shuffled
baselines of Section~\ref{ch10:sec:blinded} exist to estimate. A
priority order is not a significance ranking, and the candidate emits
no p-values here.

This is a deliberately modest job description, and the candidate keeps
it modest structurally: residual outputs feed no gate, trigger no
label, and modify no profile. They are inputs to human research
planning, Ring~4, and their provenance chains exist so that the human
doing the planning can see exactly what stands behind each number.
The asymmetry with the rest of the pipeline is intentional: everywhere
else, outputs are consumed by further typed operations; here, the
consumer is a person deciding where to measure next, and the type
system's last act is to hand that person an honest record rather than
a conclusion.

\section{The blinded implanted-fault program}\label{ch10:sec:blinded}

The registered residual-localization hypothesis
(\textsc{clm-cas-ohm-residual-localization}) is the candidate's one
sanctioned route by which residuals could ever earn more than the
triage role --- and it is constructed so that failure is meaningful.
The program, fixed in the preregistered closure protocol and the
synthetic known-fault infrastructure, has the following shape.

Faults with known ground truth are implanted into synthetic boards ---
hidden constraints, bypasses, cutoffs, saturations, off-axis couplings,
non-ohmic responses --- by a process separated from the analysis: the
analyst receives sweeps and residual diagnostics without fault labels.
The diagnostic pipeline ranks candidate fault locations from residual
evidence alone. Performance is scored against the implantation record,
and --- this is the load-bearing clause --- against two null references:
\emph{matched controls} (boards identical in construction with no
implanted fault, so that the false-positive behavior of the ranking is
measured on truly clean substrate) and \emph{shuffled baselines}
(residual-to-location assignments permuted, so that the information
content of the actual assignment is measured against chance under the
same marginals). Localization is claimed only if ranking performance
exceeds both references at the preregistered margin; negative outcomes,
exclusions, and instability are recorded without erasure, per the
closure protocol's general terms.

\begin{protocol}[Withdrawal clause]\label{ch10:prot:withdrawal}
If the blinded program fails --- ranking not above shuffled baseline, or
false-positive behavior on matched controls incompatible with the
preregistered margin --- causal language about residuals is withdrawn,
not weakened. The registered hypothesis is closed as refuted; the
non-identifying contract of Rule~\ref{ch10:rule:nonident} remains the
permanent ceiling on residual semantics; and no rephrasing (``suggests,''
``indicates,'' ``is consistent with a bypass'') survives as licensed
language, because each of those phrasings is precisely the mechanism
attribution the failed test failed to earn.
\end{protocol}

The synthetic substrate for the program already exists in the bounded
software evidence and is worth describing, because it fixes what
``implanted fault'' means operationally. A synthetic board is a typed
generator with a declared base resistance and a declared fault mode
drawn from a closed enumeration --- none, hidden constraint (base
opposition raised by a hidden component), bypass (base opposition
reduced), cutoff (zero flow below a drive threshold), saturation (flow
capped at a ceiling), off-axis (a declared fraction of flow diverted
off the measured axis), and non-ohmic (flow quadratic in drive) ---
each altering the board's dose-response law in a known, parameterized
way. Sweeps over such boards produce observation tuples
indistinguishable in type from any other sweep, which is the property
the blinding leans on: the analysis pipeline runs the same typed path
--- classification, prediction, residual --- that any admitted data
would traverse, and the fault label exists only in the implantation
record held outside the analysis. The known-fault suite currently
serves as golden-fixture verification (the pipeline detects each
implanted signature where detection is defined, and stays silent where
silence is required); the blinded \emph{scoring} run, with its matched
controls and shuffled baselines and preregistered margins, is the part
that has not been executed, and it is the part on which any upgrade of
residual semantics depends.

The two null references discipline different failure modes, and both
are necessary. Matched controls --- fault-free boards drawn from the
same construction distribution --- estimate the ranking's false-positive
behavior: a pipeline that finds ``faults'' in clean substrate at the
same rate as in faulted substrate has a ranking that reflects its own
noise, whatever its hit rate. Shuffled baselines --- true residuals
permuted against locations --- estimate how much of the ranking's
success is carried by the residual-to-location \emph{assignment} as
opposed to the marginal distribution of residual magnitudes: a
pipeline that scores no better than its own shuffle is exploiting
magnitude statistics that any location would satisfy. Passing one
reference and failing the other is a fail; the preregistered claim
requires exceeding both at the stated margin, and the margin is fixed
before the run, not after.

Two boundaries on the program itself, so that this section cannot be
overread. First, the synthetic boards validate \emph{software behavior}:
they establish that the pipeline detects what it was built to detect on
substrates where ground truth is available by construction. Success on
synthetic faults licenses nothing about live models; it is a necessary
precondition for the live experiment, which carries its own
preregistration, controls, and independent recomputation requirements.
Second, even a fully successful program would license localization
claims of the preregistered form only --- ranking above baselines at
stated margins --- not the unrestricted ``residual proves hidden
resistance'' reading, which remains on the forbidden-wording list
regardless of outcome.

\section{The source's residual-overinterpretation tensions, recapped}\label{ch10:sec:tensions}

The preserved residual conflict is the clearest instance in the source
corpus of a document warning against a reading and performing it
anyway, and the candidate's design is best understood as a mechanical
enforcement of the sources' own warnings against the sources' own
practice. The tensions, at evidence-record granularity:

\emph{The Guide.} On the pages where residual analysis is introduced,
residual sign is glossed as suggesting hidden opposition (negative) or
extra amplification / hidden assistance (positive) --- mechanism
vocabulary --- on the same page as the statement that the residual alone
is insufficient to identify a cause. The warning and the violation
share a page; nothing in the Guide's procedure operationalizes the
warning.

\emph{Appendix~I.} The quantity is renamed: what the Guide treats as a
hidden resistance \emph{component} is called a residual, an explicit
step toward the non-identifying posture --- and yet the same appendix
treats a positive residual on a safety board as a bypass finding, which
is mechanism attribution in its most consequential form. The appendix
also states that saturated-stage residuals are meaningless, an
observation the candidate promotes from remark to gate: the
active-region admission requirement of
Rule~\ref{ch10:rule:admission} makes the meaningless case uncomputable
rather than inadvisable.

\emph{The addendum.} The diagnostic walk places hidden content
\emph{last} --- after dose, direction, calibration, saturation, and
measurement checks --- and the falsification row for residuals requires
independently identified causes above a shuffled-control baseline
before residual-based claims are made. This is the source corpus at its
most careful, and it is the position the candidate implements: the
blinded program of Section~\ref{ch10:sec:blinded} is the addendum's
falsification row, given hard edges.

Table~\ref{ch10:tab:tensions} aligns each preserved tension with the
design element that enforces the source's own warning against the
source's own practice.

\begin{table}[htbp]
\centering
\small
\begin{tabular}{@{}p{0.44\textwidth}p{0.50\textwidth}@{}}
\toprule
Source tension & Enforcing design element \\
\midrule
Guide: residual sign glossed as hidden opposition or hidden
assistance, on the page stating residual alone is insufficient &
Non-identifying contract (Rule~\ref{ch10:rule:nonident}); no mechanism
field exists in any residual output type \\
\addlinespace
Appendix~I: hidden resistance renamed a residual, yet positive
residual on a safety board read as a bypass &
Apparent resistance quarantined under the
\texttt{APPARENT\_SAME\_SAMPLE} role, rejected by the prediction gate;
no bypass vocabulary in any emitted quantity \\
\addlinespace
Appendix~I: saturated-stage residuals are meaningless &
Active-region admission gate; residual computation is impossible, not
inadvisable, outside \texttt{ACTIVE} \\
\addlinespace
Addendum: hidden content last in the diagnostic walk &
Ranking-only license (Section~\ref{ch10:sec:goodfor}); residuals
allocate probes, they do not name contents \\
\addlinespace
Addendum: residual validity requires independently identified causes
above a shuffled-control baseline &
Blinded implanted-fault program with matched controls, shuffled
baselines, preregistered margins, and the withdrawal clause
(Protocol~\ref{ch10:prot:withdrawal}) \\
\addlinespace
Both: ratios reported at negligible response &
Censoring at zero-compatible observation
(Definition~\ref{ch10:def:censoring}) and the residual floor on
predicted denominators \\
\bottomrule
\end{tabular}
\caption{Preserved residual-overinterpretation tensions and their
mechanical enforcement in the candidate.}
\label{ch10:tab:tensions}
\end{table}

The recapitulation matters because it locates the candidate's
contribution precisely. The insight that residuals are non-identifying
is \emph{in the sources}; what the sources lack is any mechanism that
makes the insight binding on their own procedures. The candidate's
residual layer is that mechanism: admission gates that exclude the
meaningless cases, censoring that excludes the undefined ones, typed
outputs that cannot carry mechanism labels, a method identity that
names the contract, and a preregistered program with a withdrawal
clause as the only path to anything more. Where the Guide says
``insufficient'' and proceeds, the candidate says ``insufficient'' and
stops.

\begin{boundary}[Chapter boundary]\label{ch10:bound:residuals}
No residual over live-model data has been admitted; every residual
computed to date in the candidate's history is a synthetic-board or
test-fixture artifact within the bounded software evidence. The residual
layer's guarantees are structural: well-defined subtraction under full
derivation binding, censoring at zero-compatible denominators,
propagated uncertainty everywhere, and the impossibility --- at the type
level --- of a residual that names its cause. Whether residual
diagnostics can localize implanted faults above chance is a registered,
open, falsifiable question, and this monograph asserts no answer to it.
\end{boundary}
